%% ============================================================
%%  Аналитический отчёт
%%  Токенизация в Vision-Language-Action моделях
%%  Апрель 2026
%% ============================================================
\documentclass[12pt,a4paper]{article}

\usepackage[utf8]{inputenc}
\usepackage[T2A]{fontenc}
\usepackage[russian,english]{babel}
\usepackage{geometry}
\usepackage{setspace}
\usepackage{hyperref}
\usepackage{enumitem}
\usepackage{array}
\usepackage{booktabs}
\usepackage{titlesec}
\usepackage{xcolor}
\usepackage{fancyhdr}
\usepackage{amsmath}
\usepackage{caption}
\usepackage{listings}

\lstdefinestyle{shell}{
    basicstyle=\small\ttfamily,
    breaklines=true,
    breakatwhitespace=true,
    frame=single,
    rulecolor=\color{gray!40},
    backgroundcolor=\color{gray!6},
    xleftmargin=6pt,
    xrightmargin=6pt,
    aboveskip=4pt,
    belowskip=4pt,
}

\geometry{left=25mm, right=25mm, top=30mm, bottom=30mm}
\onehalfspacing

\hypersetup{
    colorlinks=true,
    linkcolor=blue,
    urlcolor=blue,
    citecolor=blue,
    pdftitle={Токенизация в VLA и ассистенты исследователя},
}

\titleformat{\section}{\large\bfseries}{\thesection.}{0.6em}{}
\titleformat{\subsection}{\normalsize\bfseries}{\thesubsection.}{0.6em}{}
\titleformat{\subsubsection}{\normalsize\itshape\bfseries}{\thesubsubsection.}{0.6em}{}

\pagestyle{fancy}
\fancyhf{}
\rhead{\small Токенизация в VLA}
\lhead{\small Аналитический отчёт}
\cfoot{\thepage}

%% ============================================================
\title{%
    \textbf{Токенизация в Vision-Language-Action моделях}\\[0.4em]
    \large Обзор инструментов, исследовательская гипотеза и результаты симуляции
}
\author{}
\date{Апрель 2026}

%% ============================================================
\begin{document}

\selectlanguage{russian}
\maketitle
\thispagestyle{empty}

\begin{abstract}
В работе рассматриваются подходы к токенизации в Vision-Language-Action (VLA)
моделях для управления роботами. Дан сравнительный обзор четырёх ассистентов
исследователя (\texttt{Deep Research}, \texttt{Claude Code},
\texttt{autoresearch}, \texttt{AutoResearchClaw}) и проведён анализ трёх работ:
\textit{BLT}~\cite{blt}, \textit{H-Net}~\cite{hnet} и \textit{OAT}~\cite{oat}.
Из синтеза этих работ сформулирована гипотеза об \textit{adaptive action token
refinement}: адаптивный критерий остановки при декодировании OAT tokens позволяет
экономить вычисления на простых действиях без потери качества на сложных.
Гипотеза проверена в численной симуляции на 300~эпизодах трёх типов задач;
код и результаты доступны в репозитории~\cite{repo}.
\end{abstract}

\tableofcontents
\newpage

%% ============================================================
\section{Введение}

Современные VLA-модели для управления роботами используют дискретную
токенизацию action space. Открытым остаётся вопрос об эффективности:
\textit{нужно ли тратить одинаковое число вычислений на каждое действие?}
Удержание позиции и сложный манипуляционный захват интуитивно требуют
разного числа «уточнений» при декодировании.

Для изучения этого вопроса были рассмотрены три работы с разными подходами:
\texttt{BLT}~\cite{blt} переносит принцип adaptive compute allocation
на уровень байтовых патчей; \texttt{H-Net}~\cite{hnet} делает сегментацию
обучаемой end-to-end; \texttt{OAT}~\cite{oat} строит coarse-to-fine иерархию
action tokens с prefix-decodability.

Практическая часть: выбран \texttt{OAT} как базовый метод, сформулирована
гипотеза об адаптивном критерии остановки, реализована численная симуляция
для её проверки на синтетических данных.

%% ============================================================
\section{Обзор инструментов исследователя}

\subsection{Подход к оценке}

Каждый инструмент оценивался через конкретный вопрос: \textit{помогает ли он при
разборе BLT, H-Net, OAT, при формулировке гипотезы и при переходе к реализации?}
Одного универсального инструмента для всех этапов не нашлось --- для анализа
литературы и для работы с кодом нужны принципиально разные режимы.

\subsection{Разобранные инструменты}

\textbf{Deep Research}~\cite{deepresearch} наиболее полезен на этапе обзора литературы:
быстро собирает картину из нескольких работ, находит связи между ними, помогает
сформулировать гипотезу. Ограничение: не заменяет инженерную работу с репозиторием.

\textbf{Claude Code}~\cite{claudecode} --- главный инструмент кодовой части. Хорошо
живёт внутри репозитория, помогает ориентироваться в чужом коде, подготавливать
инструкции воспроизведения, локализовать точки модификации. Ограничение: не задаёт
исследовательскую рамку сам по себе.

\textbf{autoresearch}~\cite{autoresearch} интересен как концепция agentic
experimentation: короткий итерационный loop (изменить --- запустить --- измерить).
Для VLA-тематики применяется лишь косвенно --- не готовый стек под embodied-задачи.

\textbf{AutoResearchClaw}~\cite{autoresearchclaw} претендует на замыкание полного
исследовательского цикла. Как ориентир на full-cycle automation --- полезен.
Как практический инструмент для компактного проекта --- избыточен по сложности настройки.

\subsection{Итоговый рабочий стек}

\begin{table}[h!]
\centering
\caption{Сравнение инструментов по применимости}
\begin{tabular}{p{3.2cm} p{4.4cm} p{5.6cm}}
\toprule
\textbf{Инструмент} & \textbf{Где полезен} & \textbf{Ограничение} \\
\midrule
\texttt{Deep Research}    & обзор литературы, гипотеза      & не заменяет кодовую часть \\
\texttt{Claude Code}      & репозиторий, реализация         & не задаёт исследовательскую рамку \\
\texttt{autoresearch}     & agentic experimentation         & не адаптирован под VLA \\
\texttt{AutoResearchClaw} & full-cycle pipeline             & избыточен для коротких проектов \\
\bottomrule
\end{tabular}
\end{table}

Выбранная связка: \texttt{Deep~Research} для анализа статей и формулировки
гипотезы~+ \texttt{Claude Code} для инженерной части. Все ключевые решения
принимались самостоятельно, ассистенты ускоряли структурирование и оформление.

%% ============================================================
\section{Анализ статей}

\subsection{BLT: Byte Latent Transformer}

\texttt{BLT}~\cite{blt} заменяет фиксированный токенизатор динамическими патчами
переменного размера из сырых байтов. Три компоненты: лёгкий local encoder,
тяжёлый latent transformer (работает только на уровне патчей) и лёгкий local decoder.

Размер патча определяется энтропией: предсказуемые участки объединяются в крупный
патч, сложные --- дробятся. Это концентрирует compute там, где он действительно нужен.

\textbf{Применимость к VLA.} Прямой нет --- \texttt{BLT} разработан для языкового
моделирования. Но принцип \textit{не тратить одинаковое число вычислений на все
позиции} напрямую переносится как идея для adaptive action tokenization.

\subsection{H-Net: Dynamic Chunking for Hierarchical Sequence Modeling}

\texttt{H-Net}~\cite{hnet} встраивает boundary detection в архитектуру модели и
учит его end-to-end. Сегментация становится не внешней эвристикой, а обучаемой
частью модели. Иерархический дизайн позволяет одному блоку рекурсивно
применяться на разных уровнях абстракции.

\textbf{Применимость к VLA.} Теоретически богатая, но инженерно дорогая для
embodied-задач: boundary detector нужно обучать совместно с policy, что требует
существенного переписывания pipeline.

\subsection{OAT: Ordered Action Tokenization}

\texttt{OAT}~\cite{oat} --- наиболее прикладная из трёх работ для темы управления
роботами. VQ-токенизатор обучается на демонстрационных данных и упорядочивает
токены по убыванию «информационной плотности» (coarse-to-fine). Ключевое свойство:
для декодирования достаточно любого префикса --- OAT1 даёт грубое действие,
OAT8 --- точное. Policy обучается авторегрессионно предсказывать эти токены.

\textbf{Слабое место.} Число генерируемых токенов фиксируется глобально в конфиге
при обучении. Нет механизма учёта сложности конкретного действия на inference.

\subsection{Синтез}

Все три работы атакуют одну проблему с разных сторон.
\texttt{BLT} показывает, что фиксированная дискретизация ограничивает
вычислительную эффективность;
\texttt{H-Net} --- что сегментация может и должна быть обучаемой;
\texttt{OAT} --- что action tokens следует упорядочивать и делать prefix-decodable.
Именно такое пересечение идей сделало возможным формулировку гипотезы,
объединяющей все три направления.

%% ============================================================
\section{Исследовательская гипотеза}

\subsection{Формулировка}

Из синтеза трёх работ следует:

\begin{quote}
Если вместо фиксированного числа action tokens в \texttt{OAT} использовать
адаптивный критерий остановки, зависящий от сложности текущего действия,
то на простых действиях можно значительно сократить число генерируемых токенов
при приемлемой потере точности реконструкции.
\end{quote}

\subsection{Критерий остановки}

Предложен абсолютный порог, масштабируемый к величине действия:
\begin{equation}
    \text{stop when} \quad \hat{e}_k < \varepsilon \cdot \|\mathbf{a}\|_2
    \label{eq:stopping}
\end{equation}
где $\hat{e}_k$ --- ошибка реконструкции после $k$ токенов,
$\|\mathbf{a}\|_2$ --- L2-норма исходного вектора действия,
$\varepsilon$ --- настраиваемый порог.

Смысл масштабирования: простые действия имеют малую норму, поэтому порог
достигается быстрее. Сложные действия (большая норма) требуют больше токенов
для того же относительного качества. Критерий~\eqref{eq:stopping} автоматически
адаптируется к масштабу действия без отдельной настройки под каждый тип задачи.

\subsection{Связь с базовыми статьями}

\begin{itemize}[leftmargin=1.4cm]
    \item \texttt{BLT}: принцип adaptive compute allocation по сложности входа ---
          перенесён с уровня байтовых патчей на уровень action tokens.
    \item \texttt{OAT}: prefix-decodability уже поддерживается архитектурно;
          критерий остановки добавляется как post-hoc правило поверх inference.
    \item \texttt{H-Net}: идея адаптивной, зависящей от контекста сегментации ---
          упрощена до эвристического критерия без переобучения.
\end{itemize}

%% ============================================================
\section{Симуляция}

\subsection{Мотивация}

Полный запуск \texttt{OAT} требует GPU, MuJoCo и датасета LIBERO. Вместо
этого разработана численная симуляция, воспроизводящая ключевую динамику
VQ-кодирования: coarse-to-fine убывание ошибки, зависящее от сложности действия.
Это позволяет проверить \textit{логику} гипотезы без обращения к реальной среде.

Симуляция реализована в двух вариантах: \texttt{simulate\_adaptive\_oat.py}
(Python 3.8+, без внешних зависимостей) и \texttt{simulate\_adaptive\_oat.ps1}
(PowerShell, запускается на любом Windows без установки). Код и результаты
хранятся в репозитории~\cite{repo}.

\subsection{Модель}

\textbf{Траектории.} Генерировались синтетические траектории трёх типов:
\begin{itemize}[leftmargin=1.4cm]
    \item \textit{simple} --- удержание позиции, малые корректирующие движения
          (сложность $c \in [0.05,\, 0.25]$);
    \item \textit{complex} --- захват, перемещение, ротация
          ($c \in [0.60,\, 0.95]$);
    \item \textit{mixed} --- чередование фаз с синусоидальным профилем активности.
\end{itemize}
Каждый шаг --- 7D вектор действия: $a_i \sim \mathcal{N}(0,\; 0.08 + 0.85c)$.

\textbf{VQ-кодирование.} Моделируется геометрическое убывание ошибки реконструкции
при добавлении каждого токена. Скорость убывания (decay) зависит от сложности:
\begin{itemize}[leftmargin=1.4cm]
    \item для простых ($c \approx 0.15$): decay $\approx 0.30$--$0.45$ на токен
          (быстрая сходимость);
    \item для сложных ($c \approx 0.78$): decay $\approx 0.62$--$0.75$ на токен
          (медленная сходимость).
\end{itemize}
Это соответствует coarse-to-fine свойству OAT: первый токен даёт наибольшее
улучшение, каждый следующий --- меньше.

\textbf{Базовые методы.} OAT с фиксированным числом токенов: OAT1, OAT2, OAT4, OAT8.

\textbf{Adaptive OAT.} Применяется критерий~\eqref{eq:stopping} с $\varepsilon = 0.05$,
$k_{\min}=1$, $k_{\max}=8$.

\textbf{Параметры запуска.} seed=42, 100 эпизодов на каждый тип задачи
(300 итого), 20 шагов в эпизоде.

\subsection{Результаты основного эксперимента}

\begin{table}[h!]
\centering
\caption{Результаты при $\varepsilon = 0.05$, seed\,=\,42 (300 эпизодов)}
\label{tab:main}
\begin{tabular}{lccccccc}
\toprule
\textbf{Тип задачи} & \textbf{Сложн.} & \textbf{OAT1} & \textbf{OAT4} & \textbf{OAT8}
    & \textbf{Adapt.} & \textbf{Токены} & \textbf{Экономия} \\
\midrule
simple  & 0.15 & 0.1063 & 0.0094 & 0.0004 & 0.0179 & 3.0 / 8 & \textbf{62\%} \\
complex & 0.78 & 0.8019 & 0.2523 & 0.0562 & 0.0785 & 6.7 / 8 & 17\% \\
mixed   & 0.61 & 0.5973 & 0.1752 & 0.0372 & 0.0611 & 5.6 / 8 & 30\% \\
\midrule
\textbf{all}     & 0.51 & 0.5019 & 0.1456 & 0.0313 & 0.0525 & \textbf{5.1 / 8} & \textbf{36\%} \\
\bottomrule
\end{tabular}
\end{table}

Значения в столбцах OAT1/OAT4/OAT8/Adapt. --- средняя ошибка реконструкции (L2).
Столбец «Экономия» --- сокращение числа токенов относительно OAT8.

\textbf{Наблюдения по таблице:}
\begin{enumerate}[leftmargin=1.4cm]
    \item На \textit{simple} задачах adaptive метод использует в среднем \textbf{3 токена}
          вместо 8 --- экономия 62\%. По качеству реконструкции это соответствует
          уровню между OAT2 и OAT4.
    \item На \textit{complex} задачах метод расходует 6.7 токена --- сходимость
          медленнее из-за высокой нормы действия. Качество незначительно хуже OAT8
          ($+40\%$ относительно OAT8, но в абсолютных единицах: $0.079$ vs $0.056$).
    \item На \textit{mixed} задачах метод адаптируется к фазам: в фазах малой
          активности останавливается рано, в фазах высокой --- использует больше токенов.
\end{enumerate}

\subsection{Ablation: выбор порога $\varepsilon$}

\begin{table}[h!]
\centering
\caption{Влияние порога $\varepsilon$ на quality/efficiency tradeoff (все задачи)}
\label{tab:ablation}
\begin{tabular}{ccccc}
\toprule
$\varepsilon$ & \textbf{Токены} & \textbf{Экономия} & \textbf{Ошибка (адапт.)} & \textbf{$\Delta$ vs OAT8} \\
\midrule
0.02 & 6.10 & 24\% & 0.0361 & $+15\%$ \\
\textbf{0.05} & \textbf{5.11} & \textbf{36\%} & \textbf{0.0525} & $+68\%$ \\
0.10 & 3.80 & 53\% & 0.1042 & $+233\%$ \\
0.15 & 3.06 & 62\% & 0.1551 & $+396\%$ \\
0.20 & 2.46 & 69\% & 0.2090 & $+568\%$ \\
\bottomrule
\end{tabular}
\end{table}

Из таблицы~\ref{tab:ablation} следует, что $\varepsilon = 0.05$ обеспечивает
наиболее сбалансированный tradeoff: экономия 36\% токенов при росте ошибки
реконструкции на 68\% относительно OAT8. Важно учитывать масштаб: абсолютная
ошибка OAT8 очень мала (0.031), поэтому большой относительный процент не
означает значимого ухудшения качества управления.

При $\varepsilon = 0.02$ поведение близко к OAT8 ($+15\%$ ошибки) при
минимальной экономии (24\%). При $\varepsilon \geq 0.10$ экономия растёт,
но ошибка реконструкции приближается к уровню OAT2--OAT1.

\subsection{Воспроизводимость и артефакты}

Репозиторий содержит \textbf{готовые результаты} референсного запуска
в директории \texttt{results/}: для проверки не требуется повторно запускать
симуляцию. Файлы:

\begin{itemize}[leftmargin=1.4cm]
    \item \texttt{results/summary.json} --- агрегат по задачам + мета-параметры
          (seed, eps, дата);
    \item \texttt{results/episodes.csv} --- поэпизодные данные, 300 строк;
    \item \texttt{results/ablation\_eps.json} --- ablation по
          $\varepsilon \in \{0.02,\,0.05,\,0.10,\,0.15,\,0.20\}$.
\end{itemize}

Для воспроизведения из директории репозитория:

\begin{lstlisting}[style=shell, caption={PowerShell (без зависимостей)}]
powershell -ExecutionPolicy Bypass -File .\simulation\simulate_adaptive_oat.ps1 `
    -Seed 42 -NEpisodes 100 -NSteps 20 -Eps 0.05 -OutDir .\results
\end{lstlisting}

\begin{lstlisting}[style=shell, caption={Python 3.8+}]
python .\simulation\simulate_adaptive_oat.py ^
    --seed 42 --n_episodes 300 --n_steps 20 --eps 0.05 --out_dir .\results
\end{lstlisting}

Проверочные значения из \texttt{results/summary.json}:
\texttt{all.adaptive\_tokens\_mean = 5.11},
\texttt{all.token\_saving\_pct = 36.1},
\texttt{all.adaptive\_error = 0.05252}.
Мета-блок фиксирует параметры: \texttt{seed=42}, \texttt{eps=0.05},
\texttt{run\_date=2026-04-10}.

%% ============================================================
\section{Обсуждение}

\subsection{Что подтверждает симуляция}

Главный результат: \textbf{adaptive stopping работает в соответствии с гипотезой.}
На задачах типа \textit{simple} среднее число токенов составило 3.0 из 8.
Это прямое подтверждение центральной идеи: фиксированное число токенов для всех
действий неоправданно с точки зрения вычислений.

Не менее важен результат на \textit{complex} задачах: метод расходует 6.7 токена,
не экономя «там, где не стоит». Критерий~\eqref{eq:stopping} масштабируется к норме
действия автоматически --- без отдельной настройки под тип задачи.

\subsection{Ограничения}

\begin{enumerate}[leftmargin=1.4cm]
    \item \textbf{Синтетическая модель.} VQ-кодирование реального OAT может
          отличаться от геометрической модели decay. Сходимость реальных токенов
          не обязательно монотонна --- могут быть локальные осцилляции.
    \item \textbf{Нет GPU-запуска.} Baseline на LIBERO не запускался. Численные
          результаты нельзя напрямую сравнивать с success rate из оригинальной статьи OAT.
    \item \textbf{Критерий остановки эвристический.} В отличие от H-Net, он не
          учится end-to-end. Оптимальный $\varepsilon$ может зависеть от задачи.
\end{enumerate}

\subsection{Следующие шаги}
\label{sec:next_steps}

Реалистичный план продолжения при наличии GPU:
\begin{enumerate}[leftmargin=1.4cm]
    \item Запустить OAT8 baseline на LIBERO-10 (success rate как baseline метрика).
    \item Добавить критерий~\eqref{eq:stopping} в inference loop репозитория OAT.
    \item Сравнить success rate и среднее число токенов при $\varepsilon \in \{0.02, 0.05\}$.
    \item Проверить, сохраняется ли свойство «simple → меньше токенов» на реальных траекториях.
\end{enumerate}

%% ============================================================
\section{Личный вывод}

Разбирая три статьи, я несколько раз возвращался к одному вопросу: почему
авторы OAT не добавили адаптивную глубину декодирования, если prefix-decodability
уже заложена в архитектуре? Скорее всего потому, что success rate на LIBERO не
чувствителен к этому в их масштабе экспериментов. Но для более длинных горизонтов
планирования или быстрого inference это может стать важным.

Гипотеза получилась осмысленной: она не выдумана из воздуха, а прямо следует
из синтеза идей трёх независимых работ. Симуляция показала предсказуемое поведение
--- адаптивный метод действительно экономит токены именно там, где это оправдано.

Главное ограничение --- нет верификации на реальных роботических данных.
GPU-эксперимент на LIBERO оставляю как следующий шаг; для его постановки
вся необходимая база уже описана в разделе~\ref{sec:next_steps}.

%% ============================================================
\section{Заключение}

Работа прошла полный цикл: обзор инструментов, анализ трёх статей, формулировка
гипотезы, реализация симуляции и получение численных результатов. Ключевые итоги:

\begin{itemize}[leftmargin=1.4cm]
    \item Рабочий стек \texttt{Deep~Research\,+\,Claude Code} обеспечил чёткое
          разделение аналитической и инженерной фаз работы.
    \item Гипотеза об adaptive action token refinement теоретически обоснована и
          опирается на три независимые линии: adaptive compute (\texttt{BLT}),
          обучаемая сегментация (\texttt{H-Net}), prefix-decodability (\texttt{OAT}).
    \item Симуляция на 300 эпизодах подтвердила основную предпосылку:
          при $\varepsilon{=}0.05$ среднее число токенов составило 3.0 на простых
          задачах (экономия 62\%) и 6.7 на сложных --- метод адаптируется корректно.
    \item Ablation по $\varepsilon$ выявил рабочую зону $0.02$--$0.05$ как
          оптимальный баланс экономии и качества реконструкции.
    \item Все артефакты эксперимента (код, данные, логи) зафиксированы в
          репозитории~\cite{repo} и воспроизводимы одной командой.
\end{itemize}

%% ============================================================
\newpage
\begin{thebibliography}{99}

\bibitem{deepresearch}
OpenAI. Introducing Deep Research.
\url{https://openai.com/index/introducing-deep-research/}

\bibitem{claudecode}
Anthropic. Claude Code.
\url{https://claude.com/product/claude-code}

\bibitem{autoresearch}
Karpathy,~A. autoresearch.
\url{https://github.com/karpathy/autoresearch}

\bibitem{autoresearchclaw}
Aiming Lab. AutoResearchClaw.
\url{https://github.com/aiming-lab/AutoResearchClaw}

\bibitem{blt}
Pagnoni,~A. et~al.
Byte Latent Transformer: Patches Scale Better Than Tokens.
\textit{arXiv:2412.09871v1}, 2024.
\url{https://arxiv.org/abs/2412.09871v1}

\bibitem{hnet}
Hwang,~S., Wang,~B., Gu,~A.
Dynamic Chunking for End-to-End Hierarchical Sequence Modeling.
\textit{arXiv:2507.07955}, 2025.
\url{https://arxiv.org/abs/2507.07955}

\bibitem{oat}
Liu,~C. et~al.
OAT: Ordered Action Tokenization for Robotic Learning.
\textit{arXiv:2602.04215v2}, 2026.
\url{http://arxiv.org/abs/2602.04215} ·
\url{https://github.com/Chaoqi-LIU/oat}

\bibitem{repo}
Репозиторий проекта: код симуляции, результаты, логи ассистентов.
\url{https://github.com/f1b0nacc1-afk/adaptive-oat} (апрель 2026).

\end{thebibliography}

\end{document}
